% =====================================================================
%  Chapitre 2 — État de l'art
%  Règles (Sayadi/Buttler) :
%   - structure : du général au particulier, transitions entre parties
%   - chaque travail cité = quelques lignes + [n]
%   - temps : présent (faits établis), passé (allusions aux travaux)
% =====================================================================
\chapter{État de l'art}
\label{ch:soa}

% --- 2.1 Rappels cliniques -------------------------------------------
\section{Rappels cliniques sur la fibrillation auriculaire}
\label{sec:soa:clinique}

\subsection{Définition et impact}
La \FA{} se caractérise par une activité électrique désorganisée des oreillettes, entraînant un rythme ventriculaire irrégulier. Sa prévalence augmente avec l'âge et atteint près de 10~\% au-delà de 80~ans \cite{hindricks2021guidelines}.

\subsection{Signature électrocardiographique}
Sur l'\ECG{}, la \FA{} se traduit par trois signes principaux :
\begin{itemize}
  \item l'\textbf{absence d'onde P} ;
  \item la présence d'oscillations dites \textbf{de fibrillation} ;
  \item une \textbf{irrégularité des intervalles \RR{}}.
\end{itemize}

\begin{figure}[H]
  \centering
  \figplaceholder[0.85\linewidth][6cm]{ECG en rythme sinusal vs ECG en FA — illustration des trois signes}
  \caption{Exemples d'\ECG{} en rythme sinusal (haut) et en fibrillation auriculaire (bas). L'absence d'onde P et l'irrégularité des intervalles \RR{} sont caractéristiques de la \FA{}.}
  \label{fig:soa:ecg_fa}
\end{figure}

% --- 2.2 Signal RR ---------------------------------------------------
\section{Le signal \RR{} comme entrée monodimensionnelle}
\label{sec:soa:rr}

\subsection{Définition}
L'intervalle \RR{} désigne la durée séparant deux ondes R consécutives, c'est-à-dire deux contractions ventriculaires. La suite des intervalles \RR{} forme un \keyword{tachogramme} qui résume l'information de rythme indépendamment de la morphologie de l'\ECG{}.

\subsection{Avantages et limites}
L'analyse fondée sur les \RR{} présente plusieurs avantages :
\begin{itemize}
  \item robustesse au bruit basse fréquence et aux variations d'amplitude ;
  \item faible coût computationnel ;
  \item compatibilité avec les capteurs portables.
\end{itemize}
Elle souffre toutefois d'une limite intrinsèque : elle ignore la morphologie de l'onde P, dont l'absence est pourtant un marqueur direct de la \FA{}.

% --- 2.3 Jeux de données ---------------------------------------------
\section{Jeux de données publics}
\label{sec:soa:datasets}

\subsection{MIT-BIH Atrial Fibrillation Database (AFDB)}
La base AFDB \cite{moody1983new} contient 25~enregistrements \ECG{} ambulatoires de 10~heures, échantillonnés à 250~Hz, annotés battement par battement. Elle est devenue le \textit{benchmark} de référence pour la détection de la \FA{}.

\subsection{Long-Term AF Database (LTAFDB)}
La base LTAFDB \cite{petrutiu2007abrupt} contient 84~enregistrements de 24 à 25~heures, acquis dans des conditions cliniques différentes. Elle est utilisée dans ce travail comme jeu de test \emph{externe} pour évaluer la robustesse cross-dataset.

\begin{table}[H]
  \centering
  \caption{Caractéristiques des deux bases utilisées dans ce mémoire.}
  \label{tab:soa:datasets}
  \begin{tabular}{lcc}
    \toprule
    & \textbf{AFDB} & \textbf{LTAFDB} \\
    \midrule
    Nombre de patients          & 25 & 84 \\
    Durée par enregistrement    & 10 h & 24--25 h \\
    Fréquence d'échantillonnage & 250 Hz & 128 Hz \\
    Annotations                 & battement & battement \\
    Usage dans ce mémoire       & train / val / test & test externe \\
    \bottomrule
  \end{tabular}
\end{table}

% --- 2.4 Approches algorithmiques ------------------------------------
\section{Approches algorithmiques pour la détection de la \FA{}}
\label{sec:soa:approches}

\subsection{Méthodes statistiques}
Les premières approches ont exploité des descripteurs statistiques de la série \RR{} : coefficient de variation, racine de la moyenne des carrés des différences successives (RMSSD), entropie de Shannon, entropie multi-échelle \cite{tateno2001automatic,lake2011accurate}.

\subsection{Apprentissage automatique classique}
Avec l'essor de l'apprentissage automatique, ces descripteurs ont alimenté des classifieurs supervisés tels que les forêts aléatoires et le gradient boosting \cite{XX_baseline_gb}.

\subsection{Apprentissage profond sur \ECG{} brut}
Les architectures profondes appliquées directement au signal \ECG{} ont récemment atteint une performance comparable à celle d'un cardiologue \cite{hannun2019cardiologist}.

\subsection{Apprentissage profond sur intervalles \RR{}}
Plus pertinents pour ce mémoire, les travaux ciblant le seul signal \RR{} reposent généralement sur des architectures hybrides combinant convolutions 1D et couches récurrentes \cite{andersen2019deep}.

% --- 2.5 Synthèse & positionnement -----------------------------------
\section{Synthèse et positionnement}
\label{sec:soa:positionnement}

La littérature révèle trois éléments insuffisamment traités :
\begin{enumerate}[leftmargin=*]
  \item l'\textbf{évaluation à l'échelle du patient}, souvent remplacée par une évaluation par fenêtre, qui surestime la performance ;
  \item l'\textbf{empreinte mémoire} du modèle final, rarement rapportée ;
  \item la \textbf{généralisation cross-dataset}, peu étudiée.
\end{enumerate}
Le présent travail se positionne précisément sur ces trois axes.
